\begin{abstract}
Block--Accumulate--Accumulate (BAA) codes, introduced by Kolesnikov et al.
(CRYPTO 2026), combine a block-diagonal outer code with two random interleavers
and two prefix-sum accumulators.  We develop two
extensions of the BAA framework.

\noindent\quad First, over $\F_2$, we replace the sampled
block matrices of the original construction by a fixed, constant-length
constituent code with an exhaustively verified weight distribution.  This
eliminates uncertainty from sampling the constituent code and turns the global
distance calculation into an exact composition of its weight enumerator with
two binary accumulator
transitions.  Concrete Golay,
Reed--Muller, and BCH-subcode instantiations reach approximately $95.5\%$--
$99.4\%$ of the binary Gilbert--Varshamov (GV) distance while retaining
$O(N)$ encoding work.

\noindent\quad Second, we develop a large-field version using random monomial
interleavers.
In this regime, the ordinary expected weight enumerator poorly reflects the
probability that a sampled code contains a low-weight word, because such words
cluster in large projective families.  We instead group words by exact support
and shortening dimension, derive an exact
zero-return trace kernel for the accumulator, and charge only cancellation
constraints beyond the available projective dimension.  These ideas are
inspired in part by weighted nonbinary repeat-multiple-accumulate codes and by
recent large-field RAA analyses, but the block outer requires a new shortening-
dimension interface.
For $\operatorname{RS}[16,8,9]$ and $\operatorname{RS}[16,4,13]$, the resulting
support/trace first moments have working transitions near $0.30528$ and
$0.60441$.  We then quotient out trace multiplicity through a backward generic-
rank formulation.  It yields coefficient-independent structural containers at
approximately $0.30842$ and $0.60938$, respectively, caused by two-hop
neighborhoods concentrating in saturated constituent MDS blocks.  These values
remain far below the large-field GV benchmarks $0.49219$ and $0.74358$.  Thus the
current evidence suggests that sharper analysis can recover only a small amount
of slack: approaching full large-field GV likely requires a modified recursive
architecture rather than a different proof of scalar BAA.  All quoted
transitions are working numerical calculations, not interval-certified theorems
or converses for the typical ensemble.
\end{abstract}
